% registry: LIU2025
\methodhead{LIU2025}{Liu, Xie and Hao (2025): equivalent stress on the maximum shear amplitude plane}\label{sheet:LIU2025}

\block{Source} X.-Y.~Liu, J.~Xie, W.~Hao, \emph{J.\ Mater.\ Eng.\ Perform.} 34 (2025) 24315--24322
\cite{liu2025}, Eqs.~(5)--(13), (21)--(27); Fig.~3; Table~2.

\block{Idea and scope} On the plane of maximum shear stress amplitude, the shear and normal stress
amplitudes are combined as a von Mises-like equivalent stress, $\sqrt{\sigma_n^2+3\tau^2}$, and compared
with the uniaxial Basquin curve. Mean stresses enter through two corrections: Gerber for the mean shear
stress, Serensen (bilinear) for the mean normal stress. The normal stress amplitude on the critical
plane grows with the phase shift, which is how non-proportionality enters. Proposed for high-cycle
tension--torsion fatigue with phase shift, mean normal and mean shear stresses (30CrMnSiA steel,
7075-T651 aluminium alloy, TiAl alloy).

\block{Inputs}
\begin{center}\small
\begin{tabular}{@{}llp{100mm}@{}}
\toprule
Symbol & Unit & Meaning\\\midrule
$\Delta\tau_{\max}/2$, $\tau_{\max,m}$ & MPa & shear stress amplitude and mean on the critical plane\\
$\Delta\sigma_n/2$, $\sigma_{n,m}$ & MPa & normal stress amplitude and mean on the critical plane\\
$\Delta\tau'_{\max}/2=\tau'_a$ & MPa & mean-corrected shear amplitude, Eq.~\eqref{eq:l25-gerber}\\
$\Delta\sigma'_n/2=\sigma'_{n,a}$ & MPa & mean-corrected normal amplitude, Eq.~\eqref{eq:l25-serensen}\\
$R_n$ & -- & stress ratio of the normal stress on the plane, Eq.~\eqref{eq:l25-rn}\\
$\sigma_b=S_u$ & MPa & tensile ultimate strength (workbook \texttt{Su})\\
$\tau_b$ & MPa & torsional ultimate strength (workbook \texttt{Tu}, optional)\\
$S_{-1}$, $S_0$ & MPa & fully reversed and repeated limits of the Serensen diagram (Eqs.~23--24)\\
$\sigma_f'$, $b$ & MPa, -- & uniaxial Basquin coefficient and exponent\\
$N_f$ & cycles & fatigue life\\
$\lambda$ & -- & applied amplitude ratio $\sigma_a/\tau_a$ (Fig.~3)\\
$\mathrm{DP}$ & MPa & repository damage parameter, Eq.~\eqref{eq:l25-dp}\\
\bottomrule
\end{tabular}
\end{center}

\block{Formulation} The life model (Eq.~27) is
\begin{equation}
\sqrt{\Bigl(\frac{\Delta\sigma'_n}{2}\Bigr)^2 + 3\Bigl(\frac{\Delta\tau'_{\max}}{2}\Bigr)^2} = \sigma_f'(2N_f)^b ,
\label{eq:l25}
\end{equation}
with the Gerber correction of the mean shear stress (Eq.~25)
\begin{equation}
\frac{\Delta\tau'_{\max}}{2} = \frac{\Delta\tau_{\max}/2}{1-(\tau_{\max,m}/\tau_b)^2}
\label{eq:l25-gerber}
\end{equation}
and the Serensen correction of the mean normal stress (Eqs.~23, 24, 26)
\begin{gather}
R_n = \frac{\sigma_{n,m}-\Delta\sigma_n/2}{\sigma_{n,m}+\Delta\sigma_n/2} , \label{eq:l25-rn}\\
\frac{\Delta\sigma'_n}{2} =
\begin{cases}
\dfrac{\Delta\sigma_n}{2} + 0.414\,\sigma_{n,m}, & -1\le R_n\le0,\\[2mm]
\dfrac{\sqrt2\,\sigma_b\,\Delta\sigma_n/2}{\sigma_b+\Delta\sigma_n/2-\sigma_{n,m}}, & R_n>0 .
\end{cases}
\label{eq:l25-serensen}
\end{gather}
\emph{Critical plane}: the plane of maximum shear stress amplitude (Eq.~9); of its two orientations,
the one with the larger normal stress peak. The repository evaluates on every plane
\begin{equation}
\mathrm{DP} = \sqrt{(\sigma'_{n,a})^2 + 3(\tau'_a)^2}\quad(\mathrm{MPa}),
\label{eq:l25-dp}
\end{equation}
\texttt{LIU2025} on the plane of largest shear amplitude, \texttt{LIU2025\_ext} as the largest DP over
all planes. \emph{Ties} go to the larger DP, which on the source's planes is the larger normal stress
peak. \emph{Shear amplitude}: the shear amplitude and mean come from the longest chord of the shear
path (\texttt{longest\_chord}, MCC) by default; the global selection of \texttt{docs/shear-amplitude.md}
can switch the amplitude to \texttt{rss} (MCE). Without mean stresses, reversed tension gives
$\mathrm{DP}=\sigma_a$ and torsion $\mathrm{DP}=\sqrt3\,\tau_a$: the criterion implies
$\tm/\fm=1/\sqrt3$.

\block{Constants and calibration} No fitted criterion constant. Paper: the life uses the uniaxial
Basquin constants $\sigma_f'$, $b$; $\sigma_b$, $\tau_b$ from static tests (Table~1). Repository:
$\sigma_b$ from \texttt{Su}, $\tau_b$ from \texttt{Tu}, else $S_u/\sqrt3$ with a warning; the life
comes from the master curve fitted on the uniaxial $R=-1$ design tests.

\block{Life} Paper: Eq.~\eqref{eq:l25} with the uniaxial Basquin curve. Repository: the master curve
$\mathrm{DP}=AN_f^b$ (Section~\ref{sec:mastercurve}) replaces $\sigma_f'$ and $b$.

\block{Repository choices}
\begin{itemize}
\item Shear amplitude and mean from the longest chord, the paper's half range of the scalar plane shear
  (Eqs.~10--11) generalised to a shear vector.
\item $0.414$ is taken as $\sqrt2-1$, implied by Eqs.~(23)--(24) with $S_0=\sqrt2S_{-1}$ (the paper
  prints it rounded); the two branches of Eq.~(26) then join at $R_n=0$.
\item Compressive mean normal stress: Eq.~(26) does not cover $-\Delta\sigma_n/2<\sigma_{n,m}<0$, and
  its hyperbolic branch would make a fully compressive cycle damaging. The linear branch is used for
  every $\sigma_{n,m}<0$ (continuous at $\sigma_{n,m}=0$) and the corrected amplitude is clipped at 0.
\item $\tau_b$ from \texttt{Tu}; otherwise $S_u/\sqrt3$ with a warning. The estimate is coarse: the
  paper's Table~1 has $\tau_b/\sigma_b=0.78$, 0.68 and 0.94.
\item A plane with $|\tau_{\max,m}|\ge\tau_b$ (Eq.~25 undefined), or whose minimum normal stress reaches
  $\sigma_b$ (the denominator of Eq.~26 vanishes), raises \texttt{LoadingScopeError}; the model is
  skipped for that material. With the $S_u/\sqrt3$ estimate this happens on 42CrMo4\_Santus and
  Ti6Al4V\_Berto.
\item \texttt{LIU2025\_ext} is a repository variant; life from the shared master curve.
\end{itemize}

\block{Worked example} 30CrMnSiA, Table~1: $\sigma_b=1334$, $\tau_b=1043\MPa$; the
$\sigma_a=\tau_a=350\MPa$ tension--torsion tests of Fig.~5, with
$\sigma=\sigma_m+\sigma_a\sin\omega t$, $\tau=\tau_m+\tau_a\sin(\omega t-\varphi)$:
\begin{center}\small
\begin{tabular}{@{}lc@{}}
\toprule
Load & DP (MPa)\\\midrule
in phase & 700.0\\
$90^\circ$ out of phase & 700.0\\
in phase, $\tau_m=100\MPa$ & 709.7\\
in phase, $\sigma_m=100\MPa$ & 711.1\\
\bottomrule
\end{tabular}
\end{center}
In phase a larger shear amplitude acts with a smaller normal one; out of phase the reverse; the two
effects compensate exactly at $\sigma_a=\tau_a$. Both mean stresses raise DP. Reversed tension 300 and
torsion $200\MPa$ give $\mathrm{DP}=300$ and $200\sqrt3\MPa$.

\block{Verification} Fig.~3, the shear and normal amplitudes on the critical plane against the phase
angle, with each phase scaled to a largest von Mises stress of $160\MPa$ (Fig.~2): for $\lambda=0.866$
both amplitudes at all five phase angles within $2\MPa$; for $\lambda=\sqrt3$ the shear amplitude at
every phase and the normal amplitude at $0^\circ$ and $90^\circ$. At $30^\circ$ and $60^\circ$
Eqs.~(5)--(9) give 69.0 and $102.7\MPa$ against about 75 and $108\MPa$ read from the figure, also on
planes normal to the surface only: strict xfails, unresolved. Also checked: the branches of Eq.~(26)
join at $R_n=0$, Eq.~(24), the compressive-mean extension, the Gerber correction, the von Mises-like
reversed values, the missing-\texttt{Tu} warning and the scope errors. Not reproduced: the life
comparisons of Figs.~5--7 and the scatter-band fractions of Table~2 (data in the cited sources).
